\documentclass[11pt,a4paper]{article}
\usepackage[T1]{fontenc}
\usepackage[utf8]{inputenc}
\usepackage{amsmath,amssymb,xcolor}
\usepackage[margin=25mm]{geometry}
\usepackage[colorlinks=true,urlcolor=blue]{hyperref}
\definecolor{nexorablue}{HTML}{2459B5}
\setlength{\parindent}{0pt}
\setlength{\parskip}{7pt}
\setlength{\emergencystretch}{2em}
\renewcommand{\familydefault}{\sfdefault}
\begin{document}

{\small\bfseries\color{nexorablue} NEXORA CPGE}\par
{\small MPSI / NEX-MES-03 / Énoncé}\par
{\LARGE\bfseries Mesurer une résistance par quotient}\par
\vspace{0.25em}\hrule\vspace{0.6em}

On estime une résistance par $R=U/I$. Les mesures indépendantes donnent $U=5{,}00\,\mathrm V$ avec $u(U)=0{,}04\,\mathrm V$, et $I=20{,}0\,\mathrm{mA}$ avec $u(I)=0{,}2\,\mathrm{mA}$. Calculer $R$, son incertitude et la contribution dominante.


\end{document}
